% V23 verified scope: French physical90, section15 heading, through physical93
% "les résultats correspondant à 2.1. à 2.11."; English physical59--61 same boundary.
% Mathematical supplements and corrections are explicitly attributed in notes/report.
\egasection{§15\quad 前述若干结果的公理化}

我主要想到的是第2至第8节的结果：其中至少大多数，在比射影空间的超平面族（或给定次数的超曲面族）更一般的条件下仍然有效。我觉得，从正式成稿的最初几节起，就采取这样的公理化观点是适宜的。我目前还没有弄清，是否能够在这个意义下推广Bertini--Zariski（因而也推广第4节和第6节的结果）；我已经写信给这方面的专家（Serre、Zariski），询问他们是否知道这样的推广。我确实觉得，下文所考虑的那一类简单的微分性假设，应当足以推出Bertini--Zariski。如果这些专家不能给我们满意的答复，就应当尝试用我们自己的方法把问题弄清楚。\cmt{V23-SET001}{据法90末至91页恢复。英59页漏去“从正式成稿的最初几节起”，并把法文\textit{ci-dessous}（下文）译成“above”。英稿围绕“专家”的“competent people Tr”“experts”等是译编措辞痕迹，不是法稿中的另一项数学条件。法91明确写第4、6节，而本节结语写第3、6节；两处编号照录，并在报告说明差别，不擅自统一。“我”及写信的叙述均属于原作者；此段是研究计划，不是已完成的推广证明。}

我们从一个由有限表现态射组成的交换图式出发：\cmt{V23-SET002}{法91图中的右上对象为带下划线的$P$，英59的下划线遗失。本译文统一用$\mathcal P$、下文用$\mathcal X,\mathcal Z$表示这些普遍对象，并保留原图标记$(D)$。箭头名$q,p$及下文$q_X,p_X$是为核查方便补充的记号；原图没有给箭头命名。这里仅假定交换，不能据方形的外观认定它是笛卡尔图式。}
\[
(D)\qquad
\begin{tikzcd}
P \arrow[d] & \mathcal P \arrow[l,"q"'] \arrow[d,"p"] \\
S & G \arrow[l]
\end{tikzcd}
\]
（在主要应用中，$P$是射影纤维空间，$G$是由它导出的Grassmann空间，$\mathcal P$是关联预概形。\ednote{英编稿正文在\textit{grassmannienne}之后插有“形容词！”的判断，脚注68写：“可能是Grassmann纤维空间。”法91原文为\textit{une grassmannienne induite}，这里该词作名词使用；本文按法稿译出，保留上述英编意见于此。}
在最重要的情形中，相应态射
\[
\mathcal P\longrightarrow P\times_S G
\]
将是闭浸入，我们把$G$看作$P$在$S$上各纤维内的一族闭子预概形的参数概形。更确切地说，若$\xi\in G$，则$\mathcal P_\xi$是
$P_s\times_{k(s)}k(\xi)$的闭子预概形，其中$s$是$S$中位于$\xi$下方的点。\cmt{V23-SET003}{法91写\textit{en dessous de}，即$s$是$\xi$经$G\to S$的像；英60误作“over”。法稿的$P_{s,k(\xi)}$按其本义展开为剩余域扩张后的纤维。此处方向和基域扩张都有原文依据，不是凭空补加一个纤维假设。}
此外，本节大多数陈述大概都会在$S=\Spec(k)$的情形下讨论。）

在一般情形中，仍可把$G$看作$P$在$S$上各纤维之上的一族预概形的参数概形，$\xi$对应于$P_s\times_{k(s)}k(\xi)$之上的$\mathcal P_\xi$。当然，也可以不把$\xi$取为$G$的一个（绝对）点，而取为一个以$S$预概形$T$为值的点；这样便得到一个$T$态射
\[
\mathcal P_\xi\longrightarrow P_T,
\]
在最先考虑的情形中，它是闭浸入。\cmt{V23-SET004}{法92首行公式的目标印作$B_T$，英60亦如此，但本图式没有定义$B$。沿$\xi:T\to G$拉回$\mathcal P\to P\times_SG$，目标必为$P_T$，故正文订正。这里$T$可以是任意$S$预概形；不是只允许域的谱。}

若$f:X\to P$是一个态射，置
\[
\mathcal X=X\times_P\mathcal P,
\]
便得到形状同于前一方形的图式：\cmt{V23-SET005}{法92的下划线明确区分$X$与普遍对象；英60将两者排为同一字母。此纤维积构造对任意$f$成立，且由原交换图得到新的交换图；它本身不使所画方形成为笛卡尔图，也不在没有额外条件时使$X\to S$有限表现。后面的维数论证采用第1节的有限型范围，见V23-SET009。}
\[
\begin{tikzcd}
X \arrow[d] & \mathcal X \arrow[l,"q_X"'] \arrow[d,"p_X"] \\
S & G \arrow[l]
\end{tikzcd}
\]
于是显然，第2至第8节研究的所有问题，在刚提出的一般背景中仍然有意义；应当提炼出保证那些节中结论成立的公理化条件。\ednote{英编稿脚注69写：“法文\textit{dégager}，疏解、明确表述？？”法92的完整句子清楚可读，此处译为“提炼出”；英文正文漏去的动词已有法稿依据。}

我们将假设$P$和$G$在$S$上平坦，$G$的纤维几何不可约、维数为$N$（以便能够考虑泛点！），并假设$\mathcal P\to P$光滑，纤维几何不可约、维数为$N-m$。沿用余维参数$m\geq0$的范围，因而$0\leq m\leq N$。\cmt{V23-SET006}{这里将参数范围明写出来：几何不可约按非空的约定理解，纤维的维数$N-m$因而非负，且$q$满射；这不是由“光滑”单独推出满射。若任由$m<0$，后面的“对每个$m$维闭不可约子集”将成为空量词，不能用于推出非空性。$G$的几何不可约性本身不包含约化性，不得把$G$擅自当作几何整或光滑的预概形。}
于是$\mathcal X\to X$也具有同样的性质。\cmt{V23-SET007}{这是$q$沿$f$的基变换，光滑性以及几何不可约、固定维数的非空纤维均保持。法92为两个普遍对象加了下划线；英60没有保留，不能把两个态射误解为自映射。}
这里及后文所考虑的所有性质在$S$上的基变换下都是稳定的，特别可以应用于纤维。\ednote{英编稿正文将此处记为“字迹不清，且内容较长”，脚注70又写“后面的字迹不清”。法92插在\textit{envisagées}之后的手添词可辨为\textit{ici et plus tard}，即“这里及后文”；本版据原页恢复。}\cmt{V23-SET008}{这句话是作者对拟讨论性质的概括。上述有限表现、平坦、光滑、几何不可约纤维及其固定维数确实保持于基变换；对于后文才提出的公理及各个推广结论，仍须保留各自条件，不能把这句话当作任意性质都稳定的定理。}

首先设$S=\Spec(k)$，并在下述维数论证中沿用第1节的有限型范围，取$X$在$k$上有限型。\cmt{V23-SET009}{第1节法稿物理5页明确出现“有限型预概形$X$”；第15节在写任意$f$的形式构造后未重申这一点。此处显明论证所用范围，而不把形式构造限制为有限型。若删去该范围，取$P=G=\mathbb A^1_k$、$\mathcal P$为对角线、$N=m=1$及$X=\Spec k(t)$映到$P$的泛点，则$\dim X=0<m$，但泛截面非空；所以不能把下面的维数推理外推到任意预概形。}
令$Z$是$X$的一个$d$维闭子集，则它在$\mathcal X$中的逆像$\mathcal Z$是维数为
\[
d+(N-m)=N+d-m
\]
的闭子集。\cmt{V23-SET010}{$\mathcal Z\to Z$为光滑、满射且纤维几何不可约、维数$N-m$的态射。对$Z$逐个不可约分支作维数计算再取最大值，便得到本式；这里不需要$Z$纯维。不可约纤维的非空性已经包含在假设中，不能忽略这一点又把空纤维当作对原式的反例。详细核验见本版报告V23-A03。}
若$d<m$，则$\dim\mathcal Z<N$，所以$\mathcal Z\to G$不可能占优；若$\eta$是$G$的泛点，就有$\mathcal Z_\eta=\varnothing$。事实上，这个推理还表明：将$Z$换成$\overline{f(Z)}$，若$\dim f(Z)<m$，便有$\mathcal Z_\eta=\varnothing$。\cmt{V23-SET011}{法92清楚写$\dim f(Z)<m$；英60误写$f(X)$。此处恢复法稿，不能再将$f(X)$的错误归于法文原作者。$\dim f(Z)$按其闭包的维数理解；由于有限型条件，像可构造。}

我们希望对图式$(D)$提出一个条件，保证当$\dim f(Z)\geq m$时$\mathcal Z_\eta\ne\varnothing$。看来，这应当成为本背景的一条基本公理（在（2.2）的情形中，这一性质来自一个相当特殊的整体论证）：对$P$的每一个$m$维闭不可约子集$Z$，都有\cmt{V23-SET012}{法93原式明确为$\mathcal Z_\eta\ne\varnothing$，英60末的等号是转写错误。这里$\mathcal Z$指$Z$在$\mathcal P$中的逆像。本条是额外公理，不能由前面的平坦、光滑假设推出。法92末划去了“并且更确切地说$\dim\mathcal Z=\dim Z-m$”以及随后重复引出非空性的词句；该删式本身没有写泛纤维下标$\eta$；法93还有三行密集覆打删除，未能完整辨读，未补造其内容。由这条公理推出一般像维数至少$m$时的非空性，需要处理$f(Z)$与其闭包的区别；本版报告V23-A04给出补证。}
\[
\mathcal Z_\eta\ne\varnothing.
\]

再取$X$的一个闭子集$Z$，使得$\dim f(Z)\geq m$。于是$\mathcal Z\to G$占优。若$Z$不可约，则$\mathcal Z_\eta$的维数等于\cmt{V23-SET013}{法93在“占优”的源上添有$\eta$，英61也写作$Z_\eta\to G$。该泛纤维的像当然是泛点，但原稿随后使用的是总空间的维数；正文在此用$\mathcal Z\to G$把实际论证对象写清。法稿维数式中的下划线被英稿进一步丢失，正文保留总空间、原闭集和泛纤维三者的区别。}
\[
\dim\mathcal Z-\dim G=\dim Z-m.
\]
一般地，若$Z_i$遍历$Z$的不可约分支，则应逐分支取最大值。\cmt{V23-SET014}{原文只要求$Z$闭，没有要求不可约，而对所有这类$Z$一概写$\dim Z-m$是错误的；正文补明原等式的不可约条件，并写出一般情形的正确公式。反例：取$P=G=\mathbb P^1_k$，$\mathcal P$为对角线，$N=m=1$，$X=\mathbb P^1_k\amalg\mathbb A^2_k$，$f$在第一分支为恒等，在第二分支为常值。公理成立，$\dim X=2$、$\dim f(X)=1$，但泛截面只有一个点，维数为$0\ne2-1$。一般应对$Z$的各分支$Z_i$中满足$\dim f(Z_i)\geq m$者取$\max_i(\dim Z_i-m)$；补证和原句适用边界见V23-A05。}
\[
\dim\mathcal Z_\eta
 =\max_{\{i:\,\dim f(Z_i)\geq m\}}(\dim Z_i-m).
\]
这些性质使我们能够在当前背景下展开与（2.1）至（2.11）相应的结果；其中涉及局部性质沿泛截面传递的结论，还须保证泛截面到$X$的相应正则性条件。\cmt{V23-SET015}{法93写“2.1.\ à 2.11.”，是连续范围；英61误为仅“2.1 and 2.11”。但原稿对该范围的概括少了一个必要的限制：几何不可约的$G$可以非约化，使$\Spec k(\eta)\to G$不平坦。例如令$P=G=\mathcal P=X=\Spec(k[\epsilon]/(\epsilon^2))$，令$q,p,f$均为恒等，$N=m=0$；初始条件及非空公理均满足，泛截面却是正则的$\Spec k$，而$X$不正则，不能无条件推广（2.11）的局部等价。一个充分补充是$G$在泛点处约化：此时$\Spec k(\eta)\to G$为局部化型正则态射，其基变换再与$q_X$复合，给出$\mathcal X_\eta\to X$正则。这里核验这一前提及反例，不宣称本轮已逐条重审（2.1）至（2.11）的全部旧证明。}
